\section{Шүүмж ба багийн хариу}
\subsection{Ажиллуулах заавар, горим ба хэрэглэгчийн төлөв}
II-4-10 багийн шүүмжид нийт 11 зүйл бүртгэгдсэн. Үүнээс Critical түвшний асуудал 0, Major 1, Minor 5, Suggestion 5 байв. Доорх хоёр хүснэгтэд эх баримтын дугаарлалтыг хадгалж, багийн шийдвэрийг хэрэгжсэн өөрчлөлт ба түүнийг шалгасан нотолгоотой холбов.

\tbl{Reviewer Comments and Responses — асуудал 1–6}{tab:rev1}{
\begin{tabularx}{\linewidth}{@{}c L{31mm}l L{21mm}Y@{}}\toprule
\hd № & \textbf{Шүүмж} & \textbf{Түвшин} & \textbf{Шийдвэр} & \textbf{Өөрчлөлт ба нотолгоо}\\\midrule
1 & PowerShell дээр README-ийн команд файлыг уншихын өмнө ParserError өгсөн. & Major & Зөвшөөрсөн & CLI одоо \code{--commands} аргументаар файл уншина. Файл ба stdin-ийн гаралт байт бүрээр таарсан. README, Manual-ийн Жишээ 4 ба 2.5-д ижил командыг оруулав.\\\midrule
2 & Номын жишээний товч өмнөх Banker-гүй горимыг хадгалсан. & Minor & Зөвшөөрсөн & Товч, цэс өгөгдөл ба горимыг хамт ачаална. P1-д олгож, P0-д татгалзсаныг тестээр, горим сэргэснийг хөтчөөр шалгав. \gitlink{48e935d}.\\\midrule
3 & Хүснэгтийг зассан ч өмнөх төлвөөр хүсэлт илгээж болдог. & Minor & Зөвшөөрсөн & «Шалгах» хүртэл Request, Release ба горим солих үйлдэл түгжинэ. R0-ийг 3-аас 4 болгож товчнууд хэзээ нээгдэхийг хөтчөөр шалгав. \gitlink{48e935d}.\\\midrule
4 & Графын сан, туршсан орчин, дотоод баримтын мэдээлэл зөрсөн. & Minor & Зөвшөөрсөн & Графыг inline SVG гэж тодорхойлов. Manual 2.2-т Windows-ийн БД1 ба Debian-ийн БД2 дүнг салгав. Нөөцийн формат, шаардлагын хүснэгтийг зөв файлд байрлуулав. \gitlink{ac102e4}.\\\midrule
5 & Release-ийн алдаа Request тоолуурт ордог. Түүх хэзээ арилах нь тодорхойгүй. & Minor & Зөвшөөрсөн & Алдаатай Release ба Request түүхэнд хоёр мөр, Request-ийн алдаанд нэгээр бүртгэгдэнэ. Хуудсыг шинэчлэхэд түүх арилж, хуваарилалт үлдэнэ. Зураг~\ref{fig:log}, \gitlink{9cb02c2}.\\\midrule
6 & Сценарийн нэр нь UNSAFE төлвийг бодит түгжрэл мэт ойлгуулсан. & Minor & Зөвшөөрсөн & Нэрийг «Banker-гүй: UNSAFE төлөвт хүрэх» болгов. P0-д олгоход Work $(2,1,0)$ үлдэж дараалал олдохгүй. \gitlink{85fa092}.\\\bottomrule
\end{tabularx}}
Эхний зургаан асуудлаас №2 ба №3 нь нэг ижил эрсдэлтэй байв. Хэрэглэгчийн харж буй горим, тоо хоёр серверийн шийдвэрт хэрэглэсэн төлвөөс зөрж байжээ. Иймээс эдгээр засварыг хүснэгтийн тоо, чагт болон хүсэлтийн бодит хариуг хамтад нь ажиглаж баталгаажуулсан.
\newpage
\subsection{Дүрслэл, тест, хэмжээ ба нотолгооны саналууд}
Дараагийн таван санал нь графын уншигдах байдал, жишээний өгөгдөл болон оролтын шалгалтад хамаарна. Хүснэгтэд санал тус бүрийг ямар аргаар шалгасныг, мөн багийн оролцоог ямар баримтаар тогтоосныг тэмдэглэв.
\tbl{Reviewer Comments and Responses — санал 7–11}{tab:rev2}{
\begin{tabularx}{\linewidth}{@{}c L{31mm}l L{21mm}Y@{}}\toprule
\hd № & \textbf{Шүүмж} & \textbf{Түвшин} & \textbf{Шийдвэр} & \textbf{Өөрчлөлт ба нотолгоо}\\\midrule
7 & Анх нуугдсан Need сумыг харах аргыг тайлбарлах. & Suggestion & Зөвшөөрсөн & Процесс дээр хулгана аваачих, дарах, эсвэл бүх Need сумыг харуулах сонголтыг README, Manual 1.6-д тайлбарлав. Зураг~\ref{fig:graph}.\\\midrule
8 & 3×2 ба 10×5 жишээг автоматаар шалгах. & Suggestion & Зөвшөөрсөн & UI ба хоёр тест ижил JSON уншина. 3×2 жишээнд Release-ийн дараах шилжилт, 10×5-д бүх процессыг багтаасан дарааллыг шалгав. \gitlink{85fa092}.\\\midrule
9 & Хүснэгтийн хэмжээ, POST хүсэлтийн төрөл ба гарал үүслийг шалгах. & Suggestion & Зөвшөөрсөн & 20×8 хүснэгтийг зөвшөөрч, 21×8 ба 20×9-ийг татгалзав. Тохироогүй Content-Type-д 415, өөр Origin-д 403 өгч, төлөв өөрчлөх үйлдэл дуудагдаагүй. \gitlink{6d9c4a3}.\\\midrule
10 & Хуучин кодчилолтой орчинд монгол үсэг хэвлэхэд алдаа гарах эрсдэлтэй. & Suggestion & Зөвшөөрсөн & Debian дээр cp1251 тохиргоонд «ө» үсгийг бичихэд код 2 буцсаныг давтав. Засварын дараа тест гаралтыг UTF-8-аар уншиж монгол мөрүүдийг тулгав. \gitlink{6d9c4a3}.\\\midrule
11 & Гишүүдийн хувь нэмрийг commit эсвэл PR-тай холбох. & Suggestion & Хэсэгчлэн зөвшөөрсөн & Гурван гишүүний ажилд commit холбоос нэмэв. Бусдын үүргийг БД1 тайлангийн хавсралт А-аар баримтжуулав. Хүснэгт~\ref{tab:team}.\\\bottomrule
\end{tabularx}}

\lead{Шийдвэрийн нэгтгэл.} Арван заалтыг зөвшөөрч, нэгийг хэсэгчлэн зөвшөөрсөн. №4-ийн дагуу баримтуудыг тулгаж шинэчилсэн. №11-д үлдсэн гурван гишүүний үүргийг Git түүхтэй тус бүрчлэн холбох нотолгоо дутуу байна.
\newpage
